\documentclass[11pt,a4paper]{article}
\usepackage{fontspec}
\usepackage{kotex}
\usepackage{geometry}
\usepackage{graphicx}
\usepackage{xcolor}
\usepackage{tcolorbox}
\usepackage{booktabs}
\usepackage{array}
\usepackage{longtable}
\usepackage{colortbl}
\usepackage{fancyhdr}
\usepackage{titlesec}
\usepackage{hyperref}
\usepackage{emoji}
\usepackage{amsmath}
\usepackage{amssymb}
\usepackage{enumitem}

\geometry{left=25mm, right=25mm, top=30mm, bottom=30mm}

\setmainfont{Noto Sans CJK KR}
\setsansfont{Noto Sans CJK KR}
\setmonofont{Noto Sans CJK KR}
\setemojifont{Noto Color Emoji}

\definecolor{primary}{HTML}{1976D2}
\definecolor{darkbrown}{HTML}{3E2723}
\definecolor{mediumbrown}{HTML}{5D4037}
\definecolor{lightbrown}{HTML}{8D6E63}
\definecolor{cream}{HTML}{FFF3E0}
\definecolor{lightgray}{HTML}{F5F5F5}
\definecolor{tableheader}{HTML}{E8E8E8}

\titleformat{\section}
  {\Large\bfseries\color{primary}}
  {\thesection.}{0.5em}{}
\titleformat{\subsection}
  {\large\bfseries\color{darkbrown}}
  {\thesubsection}{0.5em}{}
\titleformat{\subsubsection}
  {\normalsize\bfseries\color{mediumbrown}}
  {\thesubsubsection}{0.5em}{}

\renewcommand{\contentsname}{목차}
\renewcommand{\figurename}{그림}
\renewcommand{\tablename}{표}

\tcbset{
  tipbox/.style={colback=cream,colframe=primary,boxrule=1.5pt,arc=3mm,
    left=5mm,right=5mm,top=3mm,bottom=3mm,fonttitle=\bfseries\color{primary}},
  formulabox/.style={colback=lightgray,colframe=lightgray,boxrule=0pt,arc=2mm,
    left=5mm,right=5mm,top=3mm,bottom=3mm},
  summarybox/.style={colback=white,colframe=primary,boxrule=2pt,arc=4mm,
    left=5mm,right=5mm,top=5mm,bottom=5mm}
}

\hypersetup{colorlinks=true,linkcolor=primary,urlcolor=primary,bookmarks=true,unicode=true}


\pagestyle{fancy}
\fancyhf{}
\fancyhead[R]{\small\color{lightbrown}수업 집중도 분석 프로그램 사용설명서 v1.0.0.0}
\fancyfoot[C]{\thepage}
\renewcommand{\headrulewidth}{0.4pt}

\begin{document}

\begin{titlepage}
\centering
\vspace*{3cm}
{\Huge\bfseries\color{primary}\emoji{bar-chart} 수업 집중도 분석 프로그램}\\[0.5cm]
{\Huge\bfseries\color{primary} 사용설명서}
\\[1.5cm]
{\Large 시뮬레이터 버전: v2.4.0.0}
\\[2cm]
\begin{tcolorbox}[summarybox]
\centering
{\large 이 사용설명서는 수업 집중도 분석 프로그램의 전체 기능을 상세히 설명합니다.}
\end{tcolorbox}
\vfill
{\large 사물인터넷과 빅데이터}\\[0.3cm]
{\large 청주교육대학교 컴퓨터교육과}
\vspace{1cm}
\end{titlepage}

\section*{1. 개요}

\subsection*{1.1 목적}

이 문서는 \textbf{수업집중도 분석 프로그램}의 전체 기능을 상세히 설명하는 사용설명서입니다. 시뮬레이터를 처음 사용하기 전에 읽거나, 특정 기능을 자세히 알고 싶을 때 참조합니다.


\subsection*{1.2 관련 문서}

\begin{center}
\renewcommand{\arraystretch}{1.4}
\begin{tabular}{|p{5.5cm}|p{7.3cm}|p{2.3cm}|}
\hline
\rowcolor{tableheader}
\textbf{문서} & \textbf{용도} & \textbf{깊이} \\
\hline
퀵가이드 & 소프트웨어 사용 \textbf{중} 빠른 참조 & \emoji{star} \\
\hline
\textbf{사용설명서 (본 문서)} & \textbf{사용 전/후 전체 기능 이해} & \textbf{\emoji{star}\emoji{star}\emoji{star}} \\
\hline
이론해설서 & 배경 개념·원리 학습 & \emoji{star}\emoji{star}\emoji{star}\emoji{star} \\
\hline
활용가이드 & 수업 적용 방법 & \emoji{star}\emoji{star}\emoji{star} \\
\hline
개발참조문서 & 코드 구조·기술 사양 & \emoji{star}\emoji{star}\emoji{star}\emoji{star}\emoji{star} \\
\hline
\end{tabular}
\end{center}

\subsection*{1.3 시스템 요구사항}

\begin{center}
\renewcommand{\arraystretch}{1.4}
\begin{tabular}{|p{2.6cm}|p{12.4cm}|}
\hline
\rowcolor{tableheader}
\textbf{항목} & \textbf{요구사항} \\
\hline
브라우저 & Chrome, Edge, Firefox (최신 버전) \\
\hline
화면 & 1024px 이상 권장 (반응형 지원) \\
\hline
인터넷 & Chart.js CDN 로딩을 위해 최초 1회 필요 \\
\hline
설치 & 불필요 (HTML 파일 더블클릭으로 실행) \\
\hline
\end{tabular}
\end{center}


\section*{2. 화면 구성}

\begin{center}
\includegraphics[width=0.95\textwidth]{images/S6_Concentration_screen.png}
\end{center}


\section*{3. 기능 상세}

\subsection*{3.1 데이터 수집}

샘플 로드, 랜덤 생성, CSV 업로드, 직접 입력 4가지 방식


\subsection*{3.2 데이터 분석}

통계 카드 4개 + 산점도·히스토그램·상관관계 차트


\subsection*{3.3 학생 모델}

슬라이더로 온도·소음 기준 설정 \textrightarrow{} if-else 규칙 기반 예측


\subsection*{3.4 교수 모델}

최소자승법 다중 회귀 \textrightarrow{} 수식·계수·R² 자동 계산


\subsection*{3.5 모델 비교}

학생/교수 모델 예측 곡선 비교 + 정확도 대조


\subsection*{3.6 예측 테스트}

온도·소음 입력 \textrightarrow{} 두 모델의 예측 결과 비교


\subsection*{3.7 CSV 내보내기}

생성된 데이터를 CSV 파일로 다운로드


\subsection*{3.8 차트 시각화}

Chart.js 기반 실시간 차트 (4종)



\section*{4. 조작 방법}

\subsection*{4.1 기본 워크플로우}

1. \textbf{데이터 준비}: \emoji{inbox-tray} 수집 탭에서 샘플/랜덤/CSV/직접 입력 중 선택

2. \textbf{데이터 분석}: \emoji{bar-chart} 분석 탭에서 통계·차트 관찰

3. \textbf{학생 모델}: \emoji{man-student} 탭에서 슬라이더로 규칙 설정 \textrightarrow{} 예측 테스트

4. \textbf{교수 모델}: \emoji{man-teacher} 탭에서 학습 시작 \textrightarrow{} 수식·R² 확인

5. \textbf{비교}: \emoji{balance-scale} 탭에서 두 모델 정확도 대조



\section*{5. 자주 묻는 질문 (FAQ)}

\textbf{Q: 데이터를 얼마나 넣어야 하나요?}


A: 의미 있는 분석을 위해 최소 10건 이상을 권장합니다. 샘플 데이터(12건)로 시작해보세요.


\textbf{Q: R²가 낮으면 모델이 나쁜 건가요?}


A: 반드시 그렇지 않습니다. 실세계 데이터는 다양한 요인이 작용하므로 R²가 낮을 수 있습니다. 중요한 것은 모델의 한계를 이해하는 것입니다.


\textbf{Q: 학생 모델이 교수 모델보다 정확할 수도 있나요?}


A: 네, 데이터 패턴에 따라 가능합니다. 이것이 "정답은 하나가 아니다"라는 핵심 학습 포인트입니다.



\section*{6. 참고자료}

\begin{center}
\renewcommand{\arraystretch}{1.4}
\begin{tabular}{|p{1.5cm}|p{10.5cm}|p{3.0cm}|}
\hline
\rowcolor{tableheader}
\textbf{\#} & \textbf{자료명} & \textbf{비고} \\
\hline
1 & 수업집중도분석프로그램\_퀵가이드 & 빠른 참조 \\
\hline
2 & 수업집중도분석프로그램\_개발참조 & 기술 사양 \\
\hline
3 & 수업집중도분석프로그램\_활용가이드 & 수업 적용 \\
\hline
4 & 수업집중도분석프로그램\_이론해설서 & 배경 이론 \\
\hline
\end{tabular}
\end{center}


\vfill
\begin{center}
\small\color{lightbrown}
\textit{수업 집중도 분석 프로그램 사용설명서}
\end{center}
\end{document}